\section{Off-line zeros of every Niemeier Epstein zeta function}
\label{sec:niemeier-zeros}
The cusp contribution to a rank-twenty-four theta series does more than
prevent a direct application of the Dirichlet Riemann hypothesis: it
forces infinitely many zeros far from the Epstein symmetry line.
The required analytic result is the hybrid-universality theorem of
Nakamura and Pa\'nkowski \cite[Main Theorems~1--2]{NP2016}.

Let $L$ be any even unimodular positive-definite lattice of rank $24$,
including the Leech lattice, and put
\[
 \Theta_L(\tau)=\sum_{v\in L}q^{(v,v)/2},\qquad
 N_2(L)=\#\{v\in L:(v,v)=2\},\qquad A=\frac{65520}{691}.
\]
Since $M_{12}(\mathrm{SL}_2(\mathbb Z))$ is spanned by $E_{12}$ and
$\Delta$, comparison of the constant and first Fourier coefficients gives
\begin{equation}\label{eq:niemeier-zeta-decomposition}
 \Theta_L=E_{12}+\bigl(N_2(L)-A\bigr)\Delta,\qquad
 Z_L(w)=A\zeta(w)\zeta(w-11)+\bigl(N_2(L)-A\bigr)L(\Delta,w).
\end{equation}
Here $Z_L(w)=\sum_{v\ne0}((v,v)/2)^{-w}$ initially for $\Re w>12$.

\begin{theorem}[Niemeier zeros]
\label{thm:niemeier-offline}
For every even unimodular positive-definite rank-$24$ lattice $L$ and every
$23/2<\alpha<\beta<12$, there are positive constants $c_1,c_2,T_0$ such that
for $T\geq T_0$,
\[
 c_1T\leq\#\{w:Z_L(w)=0,\ \alpha<\Re w<\beta,\ 0<\Im w<T\}
       \leq c_2T.
\]
The upper bound counts multiplicities; the lower bound can count distinct
simple zeros.  The same assertion holds in the reflected strip
$12-\beta<\Re w<12-\alpha$.
In particular the Leech Epstein zeta function has infinitely many zeros
off $\Re w=6$, and so does every Niemeier Epstein zeta function.
No Riemann hypothesis is an assumption of this theorem.
\end{theorem}
\begin{proof}
Set $z=w-11$, $B=N_2(L)-A$, and consider the polynomial in $X$
\begin{equation}\label{eq:niemeier-hybrid-polynomial}
 P_z(X)=A\zeta(z+11)X+B L(\Delta,z+11).
\end{equation}
Both coefficients are Dirichlet series absolutely convergent for
$\Re z>1/2$.  For the second coefficient even the elementary bound
$\tau(n)=O(n^6)$ suffices: boundedness of
$y^6|\Delta(x+iy)|$ on the modular fundamental domain, and hence on the
upper half-plane, combined with Fourier inversion at $y=1/n$ proves that
bound.  Therefore $L(\Delta,z+11)$ converges absolutely already for
$\Re z>-4$.  The first coefficient is even more plainly absolutely
convergent in the required half-plane.

Neither coefficient is identically zero.  Indeed $A\ne0$, the two
$L$-series tend to $1$ as the real part tends to infinity, and
$B\ne0$ because $N_2(L)$ is an integer whereas $65520/691$ is not.
Thus $P_z$ is a nonmonomial polynomial of positive degree.
The Riemann zeta function is hybridly universal in
$1/2<\Re z<1$.  Main Theorem~1 of Nakamura--Pa\'nkowski
\cite[arXiv v4, p.~4]{NP2016} now supplies a lower bound proportional to
$T$ for the zeros of
\[
 P_z(\zeta(z))=Z_L(z+11)
\]
in every closed interior vertical strip of that region.
Their Main Theorem~2 gives the corresponding upper bound: its hypotheses
for $\zeta(z)$ are its standard meromorphic continuation, finite order,
and $\int_0^T|\zeta(\sigma+it)|^2dt=O_\sigma(T)$ for fixed
$1/2<\sigma<1$.

For clarity, the lower bound may be made a bound for simple zeros directly
from the same hybrid theorem.  Choose a small closed disc $K$ about a point
$z_0$ in the desired strip where the constant coefficient
$B L(\Delta,z+11)$ is nonzero.  Its existence follows because that
coefficient is a nonzero analytic function.  Shrink the disc so that it is
nonzero throughout $K$.  The coefficient $A\zeta(z+11)$ is nonzero there
by the Euler product.  For a sufficiently small nonzero $\epsilon$, the
function
\[
 f(z)=\frac{\epsilon(z-z_0)-B L(\Delta,z+11)}{A\zeta(z+11)}
\]
is analytic and nonvanishing on $K$, and
$P_z(f(z))=\epsilon(z-z_0)$ has exactly one simple zero inside $K$.
Theorem~E of Nakamura--Pa\'nkowski \cite[arXiv v4, p.~5]{NP2016} gives a
set of shifts $t\in[0,T]$ of measure at least $cT$ for which
$P_{z+it}(\zeta(z+it))$ uniformly approximates $P_z(f(z))$ on $K$.
Choose the error smaller than $|\epsilon|$ times the radius of $K$.
Rouch\'e's theorem then gives exactly one zero, counted with multiplicity,
in every translated disc, hence a simple zero.  One zero can account for
at most an interval of shifts of length twice that radius.  There are
therefore at least a positive multiple of $T$ distinct simple zeros.
Translating $z$ to $w=z+11$ proves the first assertion.

Finally Poisson summation gives
\[
 (2\pi)^{-w}\Gamma(w)Z_L(w)
  =(2\pi)^{-(12-w)}\Gamma(12-w)Z_L(12-w).
\]
Both gamma factors are finite and nonzero in the two strips in question.
Since the coefficients of $Z_L$ are real, reflection sends a zero $w$ in
the upper half-plane to $12-\overline w$, also in the upper half-plane,
with the same multiplicity.  This proves the reflected assertion.
\end{proof}

For the Leech lattice $N_2=0$, so the polynomial in
\eqref{eq:niemeier-hybrid-polynomial} is
$A\{\zeta(z+11)X-L(\Delta,z+11)\}$.
The theorem therefore applies to the Leech lattice without a separate
numerical search.  It locates zeros in every substrip of
$11.5<\Re w<12$ and $0<\Re w<0.5$; it does not claim that these regions
contain all zeros, or decide whether there are zeros in $\Re w>12$.
The rank-two class-number theorem of Davenport and Heilbronn cannot simply
be applied to a rank-$24$ lattice.  The proof above instead verifies the
hypotheses of a theorem that applies to the precise rank-$24$ decomposition.
